\begin{lemma}
The incident core preserves maximum degree:
$\Delta(G^{\circ})=\Delta(G)$.
\end{lemma}
\begin{proof}
Every core vertex is an original incident vertex, and
\noderef{IncidentCoreDegree} preserves its degree.  Hence each degree
appearing in the maximum for $G^{\circ}$ also appears in the maximum for
$G$, so $\Delta(G^{\circ})\leq\Delta(G)$.

For the reverse inequality, fix $v\in V(G)$.  If $d_G(v)=0$, then its degree
is bounded by the explicit value $0$ included in the maximum defining
\noderef{MaximumDegree}.  If $d_G(v)>0$, the finite incident-edge set is
nonempty; choose an identity $e$ incident with $v$.  Then $x=(v,e)$ is a
vertex of the subtype \noderef{IncidentVertex}, and
\noderef{IncidentCoreDegree} gives
$d_G(v)=d_{G^{\circ}}(x)\leq\Delta(G^{\circ})$.  Thus every value entering
the maximum for $G$, including the isolated-vertex value $0$, is at most
$\Delta(G^{\circ})$.  Therefore $\Delta(G)\leq\Delta(G^{\circ})$, proving
equality.
\end{proof}
